\documentclass[../main.tex]{subfiles}

\begin{document}
\section{Experimental Setup} \label{exp_setup}


\subsection{Research Questions} \label{research_questions}
We aim to answer the following research questions:
\begin{itemize}
    \item[(RQ1)] What is the performance of FinBERT in short sentence classification compared with the other transfer learning methods like ELMo and ULMFit?
    \item[(RQ2)] How does FinBERT compare to the state-of-the-art in financial sentiment analysis with targets discrete or continuous? 
    \item[(RQ3)] How does futher pre-training BERT on financial domain, or target corpus, affect the classification performance?
    \item[(RQ4)] What are the effects of training strategies like slanted triangular learning rates, discriminative fine-tuning and gradual unfreezing on classification performance? Do they prevent catastrophic forgetting? 
    \item[(RQ5)] Which encoder layer performs best (or worse) for sentence classification?
    \item[(RQ6)] How much fine-tuning is enough? That is, after pre-training, how many layers should be fine-tuned to achieve comparable performance to fine-tuning the whole model?
\end{itemize}

\subsection{Datasets} \label{dataset}

\subsubsection{TRC2-financial} \label{trc2}
In order to further pre-train BERT, we use a financial corpus we call TRC2-financial. It is a subset of Reuters' TRC2\footnote{The corpus can be obtained for research purposes by applying here: https://trec.nist.gov/data/reuters/reuters.html}, which consists of 1.8M news articles that were published by Reuters between 2008 and 2010. We filter for some financial keywords in order to make corpus more relevant and in limits with the compute power available. The resulting corpus, TRC2-financial, includes 46,143 documents with more than 29M words and nearly 400K sentences.

\subsubsection{Financial PhraseBank}
The main sentiment analysis dataset used in this paper is Financial PhraseBank\footnote{The dataset can be found here: https://www.researchgate.net/publication/251231364 \_FinancialPhraseBank-v10} from Malo et al. 2014 \cite{Malo2014}. Financial Phrasebank consists of 4845 english sentences selected randomly from financial news found on LexisNexis database. These sentences then were annotated by 16 people with background in finance and business. The annotators were asked to give labels according to how they think the information in the sentence might affect the mentioned company stock price. The dataset also includes information regarding the agreement levels on sentences among annotators. The distribution of agreement levels and sentiment labels can be seen on table \ref{tab:phrasebank}. We set aside 20\% of all sentences as test and 20\% of the remaining as validation set. In the end, our train set includes 3101 examples. For some of the experiments, we also make use of 10-fold cross validation.

    \subfile{tables/phrasebank.tex}
    
\subsubsection{FiQA Sentiment}
FiQA \cite{Maia2018} is a dataset that was created for WWW '18 conference financial opinion mining and question answering challenge\footnote{Data can be found here: https://sites.google.com/view/fiqa/home}. We use the data for Task 1, which includes 1,174 financial news headlines and tweets with their corresponding sentiment score. Unlike Financial Phrasebank, the targets for this datasets are continuous ranging between $[-1,1]$ with $1$ being the most positive. Each example also has information regarding which financial entity is targeted in the sentence. We do 10-fold cross validation for evaluation of the model for this dataset.

\subsection{Baseline Methods} \label{baseline_methods}
For contrastive experiments, we consider baselines with three different methods: LSTM classifier with GLoVe embeddings, LSTM classifier with ELMo embeddings and ULMFit classifier. It should be noted that these baseline methods are not experimented with as thoroughly as we did with BERT. Therefore the results should not be interpreted as definitive conclusions of one method being better. 

\subsubsection{LSTM classifiers}
We implement two classifiers using bidirectional LSTM models. In both of them, a hidden size of 128 is used, with the last hidden state size being 256 due to bidirectionality. A fully connected feed-forward layer maps the last hidden state to a vector of three, representing likelihood of three labels. The difference between two models is that one uses GLoVe embeddings, while the other uses ELMo embeddings. A dropout probability of 0.3 and a learning rate of 3e-5 is used in both models. We train them until there is no improvement in validation loss for 10 epochs.

\subsubsection{ULMFit}
As it was explained in section \ref{ulmfit}, classification with ULMFit consists of three steps. The first step of pre-training a language model is already done and the pre-trained weights are released by Howard and Ruder (2018). We first further pre-train AWD-LSTM language model on TRC2-financial corpus for 3 epochs. After that, we fine-tune the model for classification on Financial PhraseBank dataset, by adding a fully-connected layer to the output of pre-trained language model. 

\subsection{Evaluation Metrics} \label{evaluation_metrics}
For evaluation of classification models, we use three metrics: Accuracy, cross entropy loss and macro F1 average. We weight cross entropy loss with square root of inverse frequency rate. For example if a label constitutes 25\% of the all examples, we weight the loss attributed to that label by 2. Macro F1 average calculates F1 scores for each of the classes and then takes the average of them. Since our data, Financial PhraseBank suffers from label imbalance (almost 60\% of all sentences are neutral), this gives another good measure of the classification performance. For evaluation of regression model, we report mean squared error and $R^2$, as these are both standard and also reported by the state-of-the-art papers for FiQA dataset.

\subsection{Implementation Details} \label{implementation_details}
For our implementation BERT, we use a dropout probability of $p=0.1$, warm-up proportion of $0.2$, maximum sequence length of $64$ tokens, a learning rate of $2e-5$ and a mini-batch size of $64$. We train the model for 6 epochs, evaluate on the validation set and choose the best one. For discriminative fine-tuning we set the discrimination rate as 0.85. We start training with only the classification layer unfrozen, after each third of a training epoch we unfreeze the next layer. An Amazon p2.xlarge EC2 instance with one NVIDIA K80 GPU, 4 vCPUs and 64 GiB of host memory is used to train the models. 



\end{document}

